
\documentclass[11pt]{article}

\usepackage{amsmath}
\usepackage{amsfonts}
\usepackage{indentfirst}

\newcommand{\universe}{\mathcal{U}}
\newcommand{\fatdot}{\,\cdot\,}

\begin{document}

\title{What Does Univalence Actually Do?}

\author{Charles J. Geyer}

\maketitle

\section{Introduction}

We try to decode Section~2.10 of the HoTT book.
This requires understanding most or all of Chapter~2 of the book.

\section{Path Lifting}

\subsection{By Non-Dependent Functions}

Lemma~2.2.1 of the HoTT book says, given $f : A \to B$ and $x, y : A$,
there is a function
\begin{equation} \label{eq:apply-function-to-path}
   \text{ap}_f : (x =_A y) \to \bigl(f(x) =_B f(y)\bigr)
\end{equation}
This function can be read as application of $f$ to a path or as path lifting.
If $p$ is a path from $x$ to $y$, then $\text{ap}_f(p)$ is a path from
$f(x)$ to $f(y)$.

It is a common ``abuse of notation'' to write $f(p)$ instead
of $\text{ap}_f(p)$.  Usually no confusion arises.  This matches the
use of the same letter for the maps of objects to objects and morphisms
to morphisms for a category theoretic functor.
But strictly speaking, it is an abuse: $f$ maps points to points, not
paths to paths.

In category theory it is not an abuse.  A functor is a morphism between
categories.  It ``acts'' functionally on both parts of the categories,
mapping objects to objects and morphisms to morphisms (and satisfies
some laws).

In HoTT it is an abuse, because we do not want to define thingummies
analogous to categories in which the objects are point and the morphisms
are paths.  And so functions cannot map between those non-existent thingummies.

In classical homotopy theory, it is not an abuse.  If $p$ is defined to
have type $I \to A$, where $I$ is the unit interval of the real line,
then $\text{ap}_f(p) = f \circ p$.

\subsection{By Dependent Functions}

Lemma~2.3.1 of the HoTT book says, given $P : A \to \universe$ and $x, y : A$,
and $p : (x =_A y)$, there is a function $p_{\textstyle *} : P(x) \to P(y)$.

Back in Section~1.12 of the HoTT book (on identity types) this function
was called \emph{indiscernability of identicals}
(the reverse of Leibniz's principle,
the principle that equals may be substituted for equals).

The HoTT book calls this lemma the Transport Lemma.
They also introduce the notation
$$
   \text{transport}^P(p) \equiv p_{\textstyle *}
$$

The HoTT book also explains a topological view of this.
Think of the type family $P : A \to \universe$ as a fibration with
base space $A$, so $P(x)$ is the fiber over $x : A$, hence
$\sum_{x : A} P(x)$ is the \emph{total space} of the fibration with
first projection
$$
   \text{pr}_1 : \sum_{x : A} P(x) \to A
$$
given by $(x, \text{foo}) \mapsto x$.
Then, given a path $p : (x =_A y)$ and a point $u : P(x)$,
there is a path in the total space from $u$ to $p_{\textstyle *}(u)$.

This path is given explicitly by Lemma~2.3.2 {in} the HoTT book,
which it calls the Path Lifting Lemma.
Suppose $P : A \to \universe$ and $x : A$ and $u : P(x)$.
Then for all $p : (x =_A y)$, we have
\begin{subequations}
\begin{equation} \label{eq:path-lifting}
   \text{lift}(u, p) : (x, u) =_{\sum_{x : A} P(x)} (y, p_{\textstyle *}(u))
\end{equation}
satisfying
\begin{equation} \label{eq:path-lifting-too}
   \text{pr}_1(\text{lift}(u, p)) = p
\end{equation}
\end{subequations}
And indeed, \eqref{eq:path-lifting} does say that $\text{lift}(u, p)$ is
a path in the total space from $(x, u)$ to $(y, p_{\textstyle *}(u))$
as advertised.  And \eqref{eq:path-lifting-too} says the path which is lifted
(???) is $p$.

What the HoTT book means by this is explained far below (but before the
next lemma).  Given a dependent function $f : \prod_{x : A} B(x)$ they
define a non-dependent function $f' : A \to \sum_{x : A} B(x)$ by
defining $f'(x) = (x, f(x))$.
Then they consider $f'(p)$, again abuse of notation because $p$ is not
a term of type $A$, 

\section{Homotopy}

Definition~2.4.1 says if $f, g : \prod_{x : A} B(x)$ then
a \emph{homotopy} from $f$ to $g$ is a dependent function of type
\begin{equation} \label{eq:homotopy}
   (f \sim g) :\equiv \prod_{x : A} \left(f(x) =_{B(x)} g(x)\right)
\end{equation}
The HoTT book (Section~2.4) notes that the type $(f \sim g)$ is different
from the identity type $(f = g)$.  If the axiom of function extensionality
(HoTT book, Section~2.9) is assumed (or the univalence axiom, which implies
function extensionality), then the type $(f \sim g)$ and $(f = g)$ are
homotopy equivalent.  But we had better not jump to that yet because
homotopy equivalence is what we are trying to understand.

\section{Revised Down to Here}

The function
$$
   \text{idtoeqv} : (A =_\universe B) \to (A \simeq B)
$$
has argument $p : A =_\universe B$ a path in $\universe$ (not a function type)
and value in $(A \simeq B)$ defined by (2.4.11) in the HoTT book, which says,
$$
   (A \simeq B) :\equiv \sum_{f : A \to B} \text{isequiv}(f)
$$
and $\text{isequiv}(f)$ is defined by (2.4.10) in the HoTT book
$$
   \text{isequiv}(f) :\equiv
   \left( \sum_{g : B \to A} (f \circ g \sim \text{id}_B) \right)
   \times
   \left( \sum_{h : B \to A} (h \circ f \sim \text{id}_A) \right)
$$
where $\sim$ is defined by Definition~2.4.1 {in} the HoTT book
$$
   (f \sim g) :\equiv \prod_{x : A} \left(f(x) =_{P(x)} g(x)\right)
$$
where $f, g : \prod_{x : A} P(x)$.


So this is all very complicated.  First $f \sim g$ is read as a
\emph{homotopy} between $f$ and $g$.
It says $f \sim g$ implies $f(x) = g(x)$ for all $x$ where $f$ and $g$
are dependent functions of the same type.

The HoTT book notes that $f \sim g$ is different from $f = g$ unless
the axiom of function extensionality is assumed, in which case they
are equivalent (but not definitionally equal).

We can use this to translate $\text{isequiv}$ to
\begin{multline*}
   \text{isequiv}(f) \equiv
   \left( \sum_{g : B \to A} \prod_{y : B} ((f \circ g)(y) =_B y) \right)
   \\ 
   \times
   \left( \sum_{h : B \to A} \prod_{x : A} ((h \circ f)(x) =_A x) \right)
\end{multline*}
and this is a cartesian product, hence a pair, the first term of which is
a $\Sigma$-type, hence a dependent pair $(g, p)$, where
$g : B \to A$ and $p : y \mapsto ((f \circ g)(y) = y)$.
That is, $g$ is a function and $p$ is a proof that $g$ is a right inverse
of $f$, and such a proof is a function that maps $y : B$ to a path from
$(f \circ g)(y)$ to $y$.  Hence
$$
   \text{isequiv}(f) \equiv \bigl( (g, p), (h, q) \bigr)
$$
where $g, h : B \to A$ and $p$ is a proof that $g$ is a right inverse of $f$
and $q$ is a proof that $h$ is a left inverse of $f$.
Thus $\text{isequiv}$ is a really complicated beast.

Now $\text{idtoeqv}$ maps a path $r : (A = B)$ to a pair $(f, \text{foo})$
where $f : A \to B$ and $\text{foo}$ is a proof that $f$ is an equivalence,
that is $\text{foo} : \text{isequiv}(f)$.

But we can say more than that.  Lemma~2.3.1 (transport) says
for $P : A \to \universe$ and $p : (x =_A y)$ there is a function
$p_{\textstyle *} : P(x) \to P(y)$.
Sometimes (the HoTT book says) when we want to be fussy and indicate $P$
we can write $p_*$ as
$$
   \text{transport}^P(p, \fatdot) : P(x) \to P(y)
$$
and from this we get Lemma~2.3.2 (path lifting), which says
for $P : A \to \universe$ assume $x : A$ and $u : P(x)$ and $p : (x =_A y)$,
then there is a function
$$
   \text{lift}(u, p) : (x, u) =_{\left(\sum_{x : A} P(x)\right)}
   (y, p_{\textstyle *}(u))
$$
such that
$$
   \text{pr}_1(\text{lift}(u, p)) = p
$$


\end{document}
